\subsection{A normalization that cannot be a finite Frobenius trace}
\begin{finding}[\fid{F26}: rational normalization versus finite Weil spectra]\label{f:normalizedtrace}
Let $b>1$ be real and $f(x)\in\C(x)$. The eventual sequence $f(b^n)$ is
a finite exponential polynomial
\[
 \sum_{j=1}^s p_j(n)\alpha_j^n,
 \qquad p_j\in\C[n],\quad \alpha_j\in\C^*,
\]
if and only if $f$ is a Laurent polynomial.

On the \emph{actual} rank-three tensor module $M$ of \fid{F21}, set
$e_i=P_i/D$, the normalized algebraic idempotents. Then
\[
 \Tr_M(e_1e_2)=24+\frac{16}{D^2},\qquad
 \Tr_M(P_1P_2)=24D^2+16.
\]
For a prime power $q>1$ and either substitution $v=q^{n/2}$ or $v=q^{-n/2}$, the first trace
is not a finite exponential polynomial in $n$, even eventually.
It therefore cannot be the alternating trace of a fixed correspondence
against $F^n$ on any bounded finite-cohomology Weil realization.
This refutes the literal arithmetic word-trace comparison described
below; it does not assert that the GCT sources require that comparison.
\end{finding}
\begin{proof}
If $f=\sum a_jx^j$ is Laurent, then $f(b^n)=\sum a_j(b^j)^n$.
Conversely a finite exponential polynomial satisfies an eventual
constant-coefficient recurrence
\[
 \sum_{j=0}^r c_j f(b^{n+j})=0,\qquad c_0\ne0.
\]
For example its annihilator is a product of powers of the shift
operators $S-\alpha_j$, with nonzero constant term. The rational
function $R(x)=\sum_{j=0}^r c_j f(b^jx)$ vanishes at infinitely many
distinct $b^n$, so $R=0$. If $f$ has a finite nonzero pole $x_0$ of
maximal absolute value, none of $f(b^jx)$ for $j\ge1$ has a pole at
$x_0$: that would require the larger-modulus pole $b^jx_0$ of $f$.
Thus only $c_0f(x)$ has a pole there, a contradiction. A rational
function with poles only at $0$ and infinity is Laurent. Eventual
evaluation is always defined, since a rational function has only
finitely many poles. The same argument applies to $f(b^{-n})$ by
replacing $f(x)$ with $f(x^{-1})$.

In the source row matrices $N(a)$ from \fid{F21},
$\Tr(e_1e_2)=a^2/D^2$. The two scalar sectors contribute $20$ and zero.
Consequently
\[
 \Tr_M(e_1e_2)
 =20+\frac{16}{D}+\frac{4(D-2)^2}{D^2}
 =24+\frac{16}{D^2}.
\]
Multiplication by $D^2$ gives the unnormalized trace. Under the stated
substitutions, $D=q^n+2+q^{-n}$ and the sequence is
\[
 h_n=24+\frac{16q^{-2n}}{(1+q^{-n})^4}.
\]
As a rational function of $x=q^n$ this is
$24+16x^2/(x+1)^4$, which has a genuine pole at $x=-1$ and is not
Laurent. The criterion just proved excludes every finite exponential
polynomial, with arbitrary complex coefficients and Jordan-polynomial
terms allowed.

Finally let $T$ be any fixed endomorphism on a finite-dimensional
cohomology space with invertible Frobenius $F$. Its trace series is
\[
 \sum_{n\ge0}\Tr(TF^n)Z^n
 =\Tr\bigl(T(1-ZF)^{-1}\bigr),
\]
a rational function, equivalently a finite exponential polynomial in
$n$. A finite alternating sum over degrees has the same property.
Commutation of $T$ with $F$ is not needed for this obstruction; for a
Weil endomorphism it holds and only ordinary power sums are needed.
This proves the finite-realization assertion.
\end{proof}

\paragraph{The exact comparison excluded.}
Suppose a proposed finite realization assigns fixed Weil endomorphisms
$E_1,E_2$ to $e_1,e_2$ and requires, for every word $w$ and every
extension degree,
\[
 \sum_i(-1)^i\Tr\bigl(w(E_1,E_2)F^n\mid H^i\bigr)
 =\left.\Tr_M(w(e_1,e_2))\right|_{v=q^{n/2}}.
\]
The word $e_1e_2$ contradicts \fid{F26}; even imposing no idempotency
on the $E_i$ would not help. At one fixed specialization
$v=q^{1/2}$ the algebraic matrices can certainly act on a constant
Weil vector space. The resulting trace against $F^n$ is different
from substituting $v=q^{n/2}$ afresh at each extension degree.

The denominator issue is already visible in the finite Tate ring:
with $\mathbb L=v^2$, $D=\mathbb L+2+\mathbb L^{-1}$ is not a unit
of $\Q[\mathbb L,\mathbb L^{-1}]$. A bounded finite virtual sum of
Tate lines cannot have class $D^{-1}$. In contrast the actual
unnormalized trace $24D^2+16$ is a positive Laurent polynomial and
has a finite graded Tate realization. Individual matrix entries can
likewise be lifted to signed graded classes. Such entrywise lifts do
not construct compatible chain maps, multiplication functors or
categorical splittings of the normalized idempotents.

Localized motive rings, stacks with unbounded cohomology and completed
categories can accommodate rational expressions; they are not excluded
by this theorem. They require an additional finite extraction, or a
trace theorem beyond \fid{F02}/\fid{F24}, for the particular non-Laurent
family above. Indeed its convergent expansion is
\[
 h_n=24+16\sum_{j\ge0}(-1)^j\binom{j+3}{3}q^{-(j+2)n}.
\]
The independent analytic proof in the audit shows genuine generating
series poles at every $q^{j+2}$; the recurrence proof above avoids any
convergence assumption. Clearing denominators before taking traces is
another possible dictionary, but its compatibility with the designated
coefficient and operations is a separate claim.

